\chapter{Implementation and Results Audit Prompts}
\label{app:prompt-record}

\section{Implementation and Experimentation}
This example reconstructs part of the preserved implementation request. It records the intended work. The configurations and run results document which experiments were completed.

\begingroup\singlespacing\footnotesize
\begin{verbatim}
can you check what is already there in the project first?
go through the code and the previous results before changing anything.
i want to compare the individual models and the fusion methods.
check the datasets and what the labels actually mean.
check if the video and sensor inputs belong to the same sample.
make sure samples from the same group dont end up on both sides of the split.
keep the experiment settings fixed before checking the test results.
use the same evaluation setup for the models being compared in each study.
check how the predictions for training the fusion model are generated.
keep the assessment samples separate from training and model selection.
run the planned experiments only if the required checks pass.
keep the existing results and save the new work separately.
save the configs, seeds, predictions and the training details.
if something fails, keep the failed attempt and explain what happened.
also check the planned missing or degraded modality conditions.
report when fusion helps, when it does not, and when the result is unclear.
dont change the experiment just to get a better score.
keep the folders clear so i can follow what was done later.
\end{verbatim}
\endgroup

\section{Fetching and Checking Results}
This example reconstructs the request to retrieve and check the experimental results and assess options for local access. The related GitHub discussion considered transfer of the files. The historical record documents no upload.

\begingroup\singlespacing\footnotesize
\begin{verbatim}
can you fetch the saved results from the experiments?
i need the csv files, figures, matrices and the technical results for my thesis.
first check which files are actually available.
keep the results connected to the correct experiment and configuration.
check the sample ids and labels before comparing the predictions.
check the reported numbers from the saved predictions where possible.
keep results from different tasks and datasets clearly identified.
explain what the individual models and fusion methods actually did.
show the comparison tables and figures that support the conclusions.
include the failures and the results where there was no clear improvement.
mention if a result is incomplete or if the required file is missing.
explain the limitations in simple words so i can use them in the thesis.
also check what i would need to keep these results on my laptop.
check which files can go to github and what needs a separate download.
dont call the upload complete before checking the required files.
\end{verbatim}
\endgroup

\section{Archiving and Version Control of Experimental Outputs}
\label{sec:results-archiving}
The author used AI assistance to plan how to archive experimental outputs from the university server locally and track them with Git. The discussion covered CSV summaries, figures, and matrices and considered GitHub as an option for repository hosting. The available records document this planning only.

\section{Provenance}
The two reconstructions draw on the preserved implementation and audit requests. The original implementation chats were unavailable, so the precise wording and authorship of individual lines cannot be verified. Section~\ref{sec:results-archiving} summarises the subsequent discussion about access to experimental outputs. Proposed prompts for future work fall outside this record.
